\documentclass{beamer}

\usepackage{lmodern}

\usetheme{Boadilla}
\usecolortheme{spruce}
\useinnertheme{circles}

\title[Inference in GP Models]{Inference in Gaussian Process Models for Political Science}
\author[Duck-Mayr]{JBrandon Duck-Mayr}
\date{December 17, 2020}

\usepackage{subcaption}

\begin{document}

\begin{frame}
    \titlepage
\end{frame}

\begin{frame}
    \frametitle{Common problems in political science}
    \pause
    \begin{itemize}[<+->]
        \item Uncertainty about functional forms
        \item Violation of conditional independence
    \end{itemize}
\end{frame}

\begin{frame}
    \frametitle{Some previous solutions}
    \begin{itemize}
        \item Uncertainty about functional forms
        \onslide<2-3>{
            \begin{itemize}
                \item Nonparametric methods, such as KRLS {\footnotesize\textcolor{gray}{(Hainmuller and Hazlett 2014)}}
                \item Strong theoretical founding to reduce uncertainty
                \item Ignoring or ad hoc or under-justified basis functions
            \end{itemize}
        }
        \item Violation of conditional independence
        \onslide<3>{
            \begin{itemize}
                \item Lagged DV
                \item Fixed or random effects specifications
                \item Standard error adjustments
            \end{itemize}
        }
    \end{itemize}
\end{frame}

\begin{frame}
    \frametitle{Newer solution: Gaussian process models}
    \pause
    Advantages:
    \begin{itemize}
        \item Principled probabilistic framework
        \item Flexible and extensible
        \item Outperforms existing solutions {\footnotesize\textcolor{gray}{(Carlson 2020)}} 
    \end{itemize}
    \pause
    Increasingly being used to study politics:
    \begin{itemize}
        \item Carlson (2020)
        \item Gill (2020)
        \item Duck-Mayr, Garnett, and Montomery (2020)
    \end{itemize}
\end{frame}

\begin{frame}
    \frametitle{Contribution: Social science-style inference with GPs}
    GP models come from ML literature; most existing derivations/routines don't give us what we want
    \pause

    \vspace{\baselineskip}
    Carlson (2020) uses sampling for all parameters via Stan
    \pause

    \vspace{\baselineskip}
    I provide:
    \begin{itemize}
        \item Closed-form solution for posterior over coefficients when we are willing to assume linearity but are worried about conditional independence
        \item Average marginal effect derivation when we are worried about functional form
    \end{itemize}
\end{frame}

\begin{frame}
    \frametitle{Primer on GP models}
    Instead of saying \( y = f \left( X \right) + \varepsilon \) for fixed \( f \), we put \textit{GP prior} over \( f \):

    \[
        f \sim \mathcal{GP} \left( \mu \left( X \right), K \left( X, X \right) \right),
    \]

    so that for any finite set of observations $\mathbf{X}$,
    \( \mathbf{f} = f \left( \mathbf{X} \right) \) has a prior distribution

    \[
        \mathbf{f} \sim \mathcal{N} \left( \mu \left( \mathbf{X} \right), K \left( \mathbf{X}, \mathbf{X} \right) \right).
    \]

    We can then get the posterior over \( \mathbf{f} \) by applying Bayes' rule and Gaussian identities.
\end{frame}

\begin{frame}
    \frametitle{Primer on GP models: Regression example}

    \begin{figure}[htb]
        \centering
        \begin{subfigure}[b]{0.45\textwidth}
            \includegraphics[width = \linewidth]{plots/gpr-example-prior.pdf}
            \caption{The GP prior}
            \label{subfig:gpr-example-prior}
        \end{subfigure}
        \begin{subfigure}[b]{0.45\textwidth}
            \includegraphics[width = \linewidth]{plots/gpr-example-posterior.pdf}
            \caption{The GP regression posterior}
            \label{subfig:gpr-example-posterior}
        \end{subfigure}
        \caption{
            An example GP prior and posterior for the function \( f \left( x \right) = 2 \sin \left( 2x \right) + x \),
            where the data were simulated as
            \( x \sim U \left( -\pi, \pi \right), y = f \left(x\right) + \varepsilon, \varepsilon \sim \mathcal{N} \left( 0, 1 \right) \).
            Simulated data points are depicted with crosses,
            the prior (posterior) mean with a solid line and the 95\% CI with a shaded region,
            and the true \( f\left(x\right) \) with a dashed line.
        }
        \label{fig:gpr-example}
    \end{figure}

\end{frame}

\begin{frame}
    \frametitle{Primer on GP models: Classification example}

    \begin{figure}[htb]
        \centering
        \begin{subfigure}[b]{0.45\textwidth}
            \includegraphics[width = \linewidth]{plots/gpc-example-prior.pdf}
            \caption{The GP prior}
            \label{subfig:gpc-example-prior}
        \end{subfigure}
        \begin{subfigure}[b]{0.45\textwidth}
            \includegraphics[width = \linewidth]{plots/gpc-example-posterior.pdf}
            \caption{Laplace approximation of posterior}
            \label{subfig:gpc-example-posterior}
        \end{subfigure}
        \caption{
            An example GP prior and posterior for the function \( f \left( x \right) = 2 \sin \left( 2x \right) + x \),
            where the data were simulated as
            \( x \sim U \left( -\pi, \pi \right), \Pr\left(y = 1\right) = \sigma\left(f\left(x\right)\right) \),
            where \( \sigma \) is the logistic function.
            Observations receiving positive labels are depicted in the rug on the top margin,
            while observations receiving negative labels are depicted in the rug on the bottom margin;
            the prior (posterior) mean with a solid line and the 95\% CI with a shaded region,
            and the true \( f\left(x\right) \) with a dashed line.
        }
        \label{fig:gpc-example}
    \end{figure}

\end{frame}

\begin{frame}
    \frametitle{Inferential quantities of interest}
    But, social scientists don't typically want posterior over \( \mathbf{f} \).
    We want mean function parameter posteriors, or distribution of average marginal effect of \(d\)th predictor in \( \mathbf{X} \).
\end{frame}

\begin{frame}
    \frametitle{Inferential QOIs: Mean function parameters}
    For regression, with common case \( \mu \left( X \right) = X\beta \),
    if we are uninterested in covariance function parameters,
    posterior on \( \beta \) has a closed form solution, given in the paper.
    
    \vspace{\baselineskip}
    If we also want uncertainty in covariance parameters, must sample;
    in software, I provide \textit{much} faster sampling routine than Stan code from Carlson (2020).
\end{frame}

\begin{frame}
    \frametitle{Inferential QOIs: Average marginal effects}

In the regression context, the SAME is defined as

\begin{equation}
    \gamma_d \triangleq \dfrac{1}{N} \sum_{i = 1}^N \dfrac{\pd f \left( \mathbf{x}_i \right) }{\pd x_{id}}.
\end{equation}

Since the derivative of a GP is also a GP,
we can also get the distribution of the derivative of \( f \).
Then the SAME is the sum of correlated normal random variates, and its distribution is easy to derive (details in paper).
\end{frame}

\begin{frame}
    \frametitle{SAME Example}

\begin{figure}[htb]
    \centering
    \begin{subfigure}[b]{0.6\textwidth}
        \includegraphics[width = \linewidth]{plots/example-ame-link.pdf}
        \caption{AME of \(x\) on \(f\left(x\right)\)}
    \end{subfigure}
    \begin{subfigure}[b]{0.3\textwidth}
        \includegraphics[width = \linewidth]{plots/example-ame-prob.pdf}
        \caption{AME of \(x\) on \(\sigma\left(f\left(x\right)\right)\)}
    \end{subfigure}
    \caption{
        Average marginal effect (AME) of \( x \) from the simulation above.
        The left panel shows the AME of \( x \) on \( f\left(x\right) \) for both regression and classification;
        the right panel shows the AME of \( x \) on the probability of a positive outcome in classification.
        The true theoretical AME is given by the dashed line in both panels.
    }
    \label{fig:example-ames}
\end{figure}
\end{frame}

\end{document}

